\documentclass[../DoAn.tex]{subfiles}
\begin{document}

\section{Đặt vấn đề}
\label{section:1.1}

Hiện nay, công tác thi đua, khen thưởng trong các đơn vị quân đội thực hiện theo Luật Thi đua, Khen thưởng số 06/2022/QH15 \cite{luat062022} cùng các nghị định, thông tư hướng dẫn. Trong khi nhiều cơ quan đã số hoá công tác này, Học viện Khoa học Quân sự (HVKHQS) vẫn dựa chủ yếu vào các tệp Microsoft Word chia sẻ qua mạng nội bộ kèm bản giấy ký xác nhận. Phương thức này bộc lộ ba nhóm hạn chế cốt lõi được trình bày dưới đây. Các hạn chế này còn được phân tích chi tiết hơn qua những vấn đề khảo sát cụ thể tại Mục~\ref{section:2.1}.

Thứ nhất, hệ thống danh hiệu có cấu trúc chuỗi và khung xét lùi theo từng năm. Để đề nghị Bằng khen của Bộ trưởng Bộ Quốc phòng (BKBTBQP), quân nhân phải duy trì Chiến sĩ thi đua cơ sở (CSTĐCS) liên tục đủ hai năm. Chiến sĩ thi đua toàn quân (CSTĐTQ) yêu cầu ba năm CSTĐCS liên tục cộng đúng một BKBTBQP trong ba năm gần nhất. Bằng khen của Thủ tướng Chính phủ (BKTTCP) yêu cầu bảy năm CSTĐCS liên tục, đúng ba BKBTBQP và đúng hai CSTĐTQ trong bảy năm gần nhất. Cả ba danh hiệu BKBTBQP, CSTĐTQ và BKTTCP đều đòi hỏi quân nhân có thêm thành tích nghiên cứu khoa học cho mỗi năm nằm trong khung xét tương ứng. Khi rà soát thủ công, cán bộ phải tra chéo nhiều tệp Word để xác định một danh hiệu cũ còn nằm trong khung xét hiện tại hay đã nằm ngoài khung. Quy trình này vừa chậm, vừa dễ bỏ sót.

Thứ hai, công tác khen thưởng tại Học viện gồm bảy nhóm với nghiệp vụ tách bạch. Đó là cá nhân hằng năm, đơn vị hằng năm, Huy chương Chiến sĩ vẻ vang (HCCSVV) theo mốc 10/15/20 năm phục vụ, Huân chương Bảo vệ Tổ quốc (HCBVTQ) dựa trên thời gian giữ chức vụ phân theo nhóm hệ số, Kỷ niệm chương vì sự nghiệp xây dựng Quân đội nhân dân Việt Nam (một lần duy nhất theo giới tính), Huy chương Quân kỳ quyết thắng (HCQKQT, 25 năm phục vụ) và thành tích nghiên cứu khoa học. Mỗi nhóm có quy tắc đầu vào, mẫu quyết định và quy trình duyệt riêng. Khi lưu trữ bằng các tệp Word rời rạc, hệ thống khó giữ được tính nhất quán giữa các nhóm.

Thứ ba, các thao tác trên Word không được ghi lại. Khi một bản ghi được chỉnh sửa hoặc xoá, hệ thống không lưu lại ai đã thực hiện, vào lúc nào và vì sao, gây khó khăn cho việc truy vết, rà soát về sau. Phân quyền chỉ dừng ở mức thư mục dùng chung, chưa phản ánh được cây tổ chức thực tế của Học viện gồm cấp Phòng, cấp Khoa và cá nhân quân nhân.

Trước thực trạng đó, Học viện cần một phần mềm chuyên dụng cho công tác quản lý khen thưởng. Phần mềm phải bao quát toàn bộ vòng đời nghiệp vụ, từ nhập dữ liệu lịch sử qua Excel, kiểm tra điều kiện chuỗi tự động, xử lý đề xuất và phê duyệt, đến xuất quyết định và lưu vết. Đồng thời, phần mềm phải tuân thủ phân quyền nhiều cấp theo đặc thù môi trường quân sự.

\section{Mục tiêu và phạm vi đề tài}
\label{section:1.2}

\textbf{Mục tiêu chính} của đồ án là xây dựng một hệ thống web hỗ trợ Phòng Chính trị Học viện Khoa học Quân sự thực hiện đầy đủ các nhóm nhiệm vụ sau:

\begin{enumerate}
    \item \textbf{Quản lý danh mục lõi}: quân nhân, đơn vị theo cấu trúc cây Cơ quan đơn vị và Đơn vị trực thuộc, chức vụ và lịch sử biến động chức vụ.
    \item \textbf{Đề xuất và phê duyệt khen thưởng}: tạo đợt đề xuất cho cả bảy nhóm danh hiệu với bộ lọc nhanh theo đơn vị/năm/giới tính/mốc phục vụ; quy trình phê duyệt nhiều bước, lưu tệp PDF quyết định kèm số quyết định.
    \item \textbf{Tự động hoá nghiệp vụ}: kiểm tra điều kiện chuỗi (BKBTBQP, CSTĐTQ, BKTTCP đối với cá nhân; BKBTBQP và BKTTCP đối với đơn vị), sinh thông điệp gợi ý tiếng Việt, nhập Excel lịch sử theo cơ chế hai bước xem trước rồi xác nhận đảm bảo nguyên tắc all-or-nothing, xuất danh sách theo nhiều tiêu chí.
    \item \textbf{Bảo mật và truy vết}: phân quyền theo bốn cấp tương ứng đặc thù tổ chức, gồm SuperAdmin quản trị hạ tầng, Admin toàn quyền nghiệp vụ tại Phòng Chính trị, Manager đề xuất trong đơn vị và User tra cứu hồ sơ cá nhân, cùng nhật ký hệ thống toàn bộ thao tác làm thay đổi dữ liệu.
    \item \textbf{Vận hành}: sao lưu cơ sở dữ liệu định kỳ và quản lý nhật ký hệ thống.
\end{enumerate}

\textbf{Phạm vi đề tài} giới hạn trong các nhiệm vụ trực tiếp phục vụ công tác xét khen thưởng. Phần mềm vận hành tập trung trên máy chủ của Học viện, truy cập qua mạng nội bộ và không xuất bản ra Internet công cộng để đáp ứng yêu cầu bảo mật quân sự. Đồ án không bao gồm các nghiệp vụ tiền lương, phụ cấp, hồ sơ y tế, đánh giá cán bộ định kỳ và các quy trình hành chính khác. Việc tích hợp với hệ thống quản lý nhân sự tổng thể của Bộ Quốc phòng cũng nằm ngoài phạm vi do vượt khỏi trọng tâm khen thưởng.

\section{Định hướng giải pháp}
\label{section:1.3}

Phần mềm xây dựng theo mô hình client--server với hai thành phần triển khai độc lập trên cùng một máy chủ vật lý hoặc tách máy.

\begin{itemize}
    \item \textbf{Frontend} dùng Next.js 14 theo App Router \cite{nextjsdocs}, viết bằng TypeScript. Lớp giao diện chính do Ant Design đảm nhiệm, Tailwind CSS bổ sung cho bố cục, còn shadcn/ui dùng cho các thành phần đặc biệt cần kiểm soát mã nguồn.
    \item \textbf{Backend} dùng Express \cite{expressdocs} trên Node.js, viết bằng TypeScript, tổ chức theo kiến trúc phân tầng tách biệt định tuyến, xử lý nghiệp vụ và truy cập dữ liệu.
    \item \textbf{Cơ sở dữ liệu} dùng PostgreSQL \cite{pgdocs} truy cập qua Prisma ORM \cite{prismadocs}. Lược đồ được mô tả trong một tệp khai báo duy nhất, di trú dữ liệu được sinh tự động.
    \item \textbf{Xác thực và phân quyền} dùng cặp JSON Web Token \cite{rfc7519} có cơ chế làm mới luân phiên; bốn vai trò nội bộ được kiểm tra trước mỗi yêu cầu nhạy cảm.
    \item \textbf{Thông báo thời gian thực} dùng Socket.IO \cite{socketiodocs} cho các sự kiện phê duyệt, từ chối, xoá đề xuất và kết quả nhập Excel khối lượng lớn.
\end{itemize}

\section{Bố cục đồ án}
\label{section:1.4}

Phần còn lại của báo cáo được tổ chức thành năm chương:

\begin{itemize}
    \item Chương \ref{chapter:Related_works} trình bày khảo sát hiện trạng, sơ đồ use case, đặc tả các use case trọng yếu và yêu cầu phi chức năng.
    \item Chương \ref{chapter:Methodology} trình bày các thành phần công nghệ được chọn cùng lý do chọn trong ngữ cảnh PM QLKT.
    \item Chương \ref{chapter:Experiment} trình bày kiến trúc phân tầng, thiết kế CSDL, sơ đồ gói, lớp và tuần tự, kết quả kiểm thử thủ công cùng hướng dẫn triển khai.
    \item Chương \ref{chapter:SolutionAndContribution} trình bày năm đóng góp nổi bật, mỗi đóng góp gồm thực trạng, giải pháp và kết quả định lượng.
    \item Chương \ref{chapter:conclusion} tổng kết và đề xuất các hướng phát triển tiếp theo.
\end{itemize}

\end{document}
